% How spreadsheets learned to recalculate. Build: make docs   (tectonic)
% Licensed CC BY 4.0. Screens are generated by tools/screens.py.
\documentclass[11pt]{article}
\usepackage{kalkylark}
\usepackage{tikz}
\usetikzlibrary{arrows.meta,positioning,calc}

\tikzset{
  cellnode/.style={draw=ink, fill=paperbox, minimum width=12mm,
    minimum height=7mm, font=\sffamily\small, rounded corners=1pt},
  dep/.style={-{Stealth[length=2.2mm]}, thick, color=ink},
  sweep/.style={-{Stealth[length=2mm]}, dashed, color=phdim},
}

\title{\sffamily\bfseries\color{ink}How spreadsheets learned to recalculate\\[4pt]
  \Large\mdseries From VisiCalc to signals}
\author{\sffamily Set Lonnert and the Kalkylark authors}
\date{\sffamily Version 0.1 · November 2026}

\begin{document}
\maketitle
\thispagestyle{empty}

\begin{abstract}
\noindent A spreadsheet looks like a grid of numbers, but underneath it is
a small machine that decides, after every keystroke, which formulas to
compute again and in what order. This is a short history of that machine,
from LANPAR (1969) and VisiCalc's single sweep (1979), through Lotus
1-2-3's natural order and Excel's self-repairing calculation chain, to the
theory that recognised spreadsheets as build systems and the ``signals'' in
today's user-interface frameworks. A final part covers the rights in
VisiCalc and why a from-scratch reconstruction such as Kalkylark is on safe
ground. It was written as background for Kalkylark, which implements the
first two stages, and several of its illustrations are screens from that
program.
\end{abstract}

\vspace{1em}
\begin{center}
\begin{tikzpicture}[x=10.2mm, y=1mm, font=\sffamily\scriptsize]
  \draw[thick, color=ink] (0,0) -- (14,0);
  \foreach \yr/\x/\lbl/\side in {
      1969/0/LANPAR/1, 1979/2.5/VisiCalc/-1, 1983/4/Lotus 1-2-3/1,
      1985/5.5/Excel (Mac)/-1, 1991/7/Improv/1, 2006/9/Google Sheets/-1,
      2007/9.9/Excel MTR/1, 2018/11.6/Build systems\\à la carte/-1,
      2020/12.5/Dynamic arrays\\LAMBDA/1, 2024/13.7/TC39 signals/-1} {
    \fill[ink] (\x,0) circle (0.9pt);
    \draw[color=phdim] (\x,0) -- (\x,{\side*6});
    \node[align=center, anchor={\ifnum\side>0 south\else north\fi}]
      at (\x,{\side*6.5}) {\textbf{\yr}\\\lbl};
  }
\end{tikzpicture}
\end{center}

\newpage
\tableofcontents
\newpage

% ======================================================================
\section{The problem}

A spreadsheet cell holds either a value you typed or a \emph{formula} that
computes one from other cells. Change one number, and every formula that
depends on it, directly or through other formulas, must be computed again.
That is the whole job. Two questions make it interesting:

\begin{enumerate}[label=\arabic*.]
\item \textbf{Which} formulas need to be computed again?
\item In \textbf{which order}, so that each formula sees up-to-date inputs?
\end{enumerate}

The cells and the references between them form a \emph{dependency graph}.
An arrow from \cell{A3} to \cell{A2} means ``\cell{A2} reads \cell{A3}''.
In the sheet below \cell{A1} is \typ{+A2*2}, \cell{A2} is \typ{+A3+1} and
\cell{A3} is the number 5:

\begin{center}
\begin{tikzpicture}[node distance=22mm]
  \node[cellnode] (a3) {A3 = 5};
  \node[cellnode, right=of a3] (a2) {A2 = A3+1};
  \node[cellnode, right=of a2] (a1) {A1 = A2*2};
  \draw[dep] (a3) -- node[above, font=\sffamily\scriptsize] {read by} (a2);
  \draw[dep] (a2) -- node[above, font=\sffamily\scriptsize] {read by} (a1);
\end{tikzpicture}
\end{center}

The answers to the two questions have changed several times over fifty
years. Each answer was shaped by the hardware of its time, and each one is
still in use somewhere.

% ======================================================================
\section{Prelude: LANPAR (1969)}

The first program to recalculate properly came ten years before the
personal computer spreadsheet. LANPAR, written by Rene Pardo and Remy
Landau for budgeting at Bell Canada and AT\&T, ran on time-sharing
mainframes. It already handled \emph{forward references}, where a formula
refers to a cell further down the sheet, by computing in \emph{natural
order}: each cell only after the cells it depends on.

Pardo and Landau filed a patent on the algorithm in 1970. It was first
rejected as an unpatentable mathematical method and finally granted in 1983
as US 4,398,249. By then the personal computer spreadsheet had been
reinvented without it. Later the patent's licensee sued Lotus and other
spreadsheet makers. As far as personal computers are concerned, natural
order disappeared for more than a decade.

% ======================================================================
\section{VisiCalc (1979): sweep the whole sheet}

Dan Bricklin and Bob Frankston wrote VisiCalc for the Apple~II: a 6502 at
1\,MHz and 32\,KB of memory that had to hold the program and the sheet
together. The program was written in assembly language and cross-assembled
on a time-sharing system. The recalculation it shipped with is as simple as
it could be:

\begin{lstlisting}
recalc():
    for each column c:              # or each row, with /GOR
        for each row r:
            if cell(r, c) holds a formula:
                cell.value = evaluate(cell.formula)
\end{lstlisting}

\begin{itemize}
\item \textbf{No dependency information at all.} A formula is kept as text
  and evaluated against whatever its neighbours hold at that moment.
\item \textbf{One pass, in a fixed order.} If \cell{A1} reads \cell{A2},
  then \cell{A1} is computed from \cell{A2}'s \emph{old} value. Users
  pressed \typ{!} until the numbers stopped changing, and learned to lay
  out sheets so that formulas refer up and to the left.
\item \textbf{Circular references are not detected.} A circle of formulas
  moves one step per recalculation. It converges or oscillates in plain
  view.
\item \textbf{Formulas are evaluated strictly left to right}: \typ{2+3*4}
  is 20. A parser without a precedence table is smaller, which comes from
  the same minimalism.
\end{itemize}

Here is the sheet from section~1 in Kalkylark, which by default
recalculates exactly this way. The sweep runs down the column (dashed). It
reaches \cell{A1} first, while \cell{A2} still holds its old value, but
the dependencies (solid) point the other way:

\begin{center}
\begin{tikzpicture}[node distance=4mm]
  \node[cellnode] (a1) {A1};
  \node[cellnode, below=of a1] (a2) {A2};
  \node[cellnode, below=of a2] (a3) {A3};
  \draw[dep] (a3.east) to[bend right=50] (a2.east);
  \draw[dep] (a2.east) to[bend right=50] (a1.east);
  \draw[sweep] ([xshift=-6mm]a1.north west) --
    node[left, font=\sffamily\scriptsize] {sweep} ([xshift=-6mm]a3.south west);
  \node[right=12mm of a2, font=\sffamily\scriptsize, color=ink, align=left]
    {A2 reads A3,\\A1 reads A2};
\end{tikzpicture}
\end{center}

\screen{stale}{Column order in Kalkylark: \cell{A1} shows 2, but
$(5+1)\times2=12$.}

The cost is proportional to the number of formulas on the sheet, every
time. On a 1979 sheet that was fine. Compared with a pencil and an eraser,
it was instant.

% ======================================================================
\section{Lotus 1-2-3 (1983): natural order}

Lotus 1-2-3 ran on the IBM PC, with far more memory, and it brought natural
order to the personal computer. Microsoft's Multiplan adopted it in its
second version. The idea is to compute each formula only after the formulas
it reads. A straightforward way to do that is a depth-first walk:

\begin{lstlisting}
recalc():
    mark every formula UNVISITED
    for each formula x: visit(x)

visit(x):
    if x is DONE:     return
    if x is VISITING: report CIRC; return      # a cycle
    mark x VISITING
    for each cell y that x reads: visit(y)
    x.value = evaluate(x.formula)
    mark x DONE
\end{lstlisting}

This is a topological sort done on the fly. Three things changed for users:

\begin{itemize}
\item \textbf{Forward references just work.} The layout of a sheet no
  longer matters.
\item \textbf{Circular references are detected.} A \textsf{CIRC} indicator
  appears, and an explicit iteration count lets deliberately circular models
  converge.
\item \textbf{Everything is still recomputed} after every change, just in
  a better order. Column and row order remained as options, for
  compatibility with VisiCalc sheets.
\end{itemize}

1-2-3 also gave operators the usual precedence, so \typ{2+3*4} became 14.

\screen{natural}{Kalkylark with \typ{/GON}: one recalculation, and
\cell{A1} is right.}

\screen{circ}{A deliberate circle (a fee that is 10\% of a total including
the fee) is flagged \textsf{CIRC}.}

% ======================================================================
\section{Excel (1985): only recompute what changed}

The first Excel, for the Macintosh, advertised \emph{intelligent recalc}:
after an edit, recompute only the cells that could possibly be affected.
That needs the graph in both directions. Each cell must know not just its
\emph{precedents}, the cells it reads, but also its \emph{dependents}, the
cells that read it:

\begin{lstlisting}
on edit(cell):
    mark_dirty(cell)

mark_dirty(x):
    if x is dirty: return
    x.dirty = true
    for each dependent d of x: mark_dirty(d)

recalc():
    for each dirty formula, in dependency order:
        evaluate it; x.dirty = false
\end{lstlisting}

\begin{center}
\begin{tikzpicture}[node distance=8mm and 14mm]
  \node[cellnode, fill=phnorm!40] (b1) {B1};
  \node[cellnode, right=of b1, fill=phnorm!40] (c1) {C1};
  \node[cellnode, right=of c1, fill=phnorm!40] (d1) {D1};
  \node[cellnode, below=of b1] (b2) {B2};
  \node[cellnode, right=of b2] (c2) {C2};
  \node[cellnode, right=of c2, fill=phnorm!40] (d2) {D2};
  \draw[dep] (b1) -- (c1);
  \draw[dep] (c1) -- (d1);
  \draw[dep] (b2) -- (c2);
  \draw[dep] (c2) -- (d2);
  \draw[dep] (c1) -- (d2);
  \node[left=4mm of b1, font=\sffamily\scriptsize, color=ink, align=right] {edited};
\end{tikzpicture}

{\small\sffamily\color{ink}After editing \cell{B1}, only the shaded cells
are recomputed. \cell{B2} and \cell{C2} are untouched.}
\end{center}

A small change in a large model now costs time in proportion to what it
touches, not to the size of the sheet. Later releases of 1-2-3 followed
with \emph{minimal recalculation} and background recalculation, which
computed while you kept typing.

% ======================================================================
\section{The modern Excel engine}

Microsoft has documented how Excel's engine works. It has three stages:

\begin{enumerate}[label=\arabic*.]
\item A \textbf{dependency tree}: for every formula, the cells and ranges
  it reads.
\item A \textbf{calculation chain}: a flat list of all formula cells in
  the order they should be computed. It is saved in the workbook file.
\item \textbf{Recalculation}: walk the chain and compute the dirty cells.
\end{enumerate}

The chain is not produced by a full topological sort. It \emph{repairs
itself}. If, while walking the chain, Excel meets a formula whose input has
not been computed yet, it moves that formula further down the chain and
carries on. Over time the chain settles into a valid order, and later
recalculations run straight through it.

Real workbooks complicate the clean model in several ways:

\begin{itemize}
\item \textbf{Volatile functions.} \typ{NOW}, \typ{TODAY} and \typ{RAND}
  change on their own. \typ{OFFSET} and \typ{INDIRECT} compute \emph{which}
  cells they read, as in \typ{INDIRECT("A" \& B1)}, so their dependencies
  are only known while they run. They are recomputed on every
  recalculation, together with everything downstream.
\item \textbf{Iterative calculation} is an opt-in setting for models that
  are circular on purpose.
\item \textbf{Multithreaded recalculation} (Excel 2007) runs independent
  parts of the chain in parallel. Functions that are not thread-safe, such
  as \typ{INDIRECT} and user code, force their part back onto one thread.
\end{itemize}

% ======================================================================
\section{A side branch: Lotus Improv (1991)}

Improv, on the NeXT computer, asked a different question: why should
formulas belong to \emph{cells} at all? Data lived in named,
multi-dimensional categories, such as years, regions and products, and you
wrote one formula, \typ{Profit = Revenue - Cost}, that applied to every
year and region. That removed a whole class of errors where a formula is
copied down a column and one copy is wrong. Improv sold poorly, but its
ideas live on in Quantrix, in pivot tables, and in modern modelling tools.

% ======================================================================
\section{Spreadsheets become a language (2018--)}

\begin{itemize}
\item \textbf{Dynamic arrays} (Excel, announced 2018): one formula can
  \emph{spill} a whole array into neighbouring cells. The dependency graph
  now tracks regions whose size is only known at run time.
\item \textbf{\texttt{LET} and \texttt{LAMBDA}} (2020 onwards): named
  intermediate values and user-defined functions in the formula language
  itself. The formula language became a complete programming language,
  without macros.
\item \textbf{Collaborative, server-side calculation} (Google Sheets,
  2006, and others): recalculation moves off the user's machine, and the
  graph is shared between people editing at the same time.
\end{itemize}

% ======================================================================
\section{Theory catches up: spreadsheets are build systems}

In 2018 Andrey Mokhov, Neil Mitchell and Simon Peyton Jones published
\emph{Build systems à la carte}. It shows that Excel, \typ{make}, Shake and
Bazel are the same kind of machine, differing in two choices: how to decide
\emph{what} to rebuild, and in \emph{which order}.

\begin{center}
\begin{tabularx}{\linewidth}{@{}lXX@{}}
\toprule
 & \textbf{what is recomputed} & \textbf{in which order} \\
\midrule
VisiCalc      & everything                 & fixed sweep, by column or row \\
Lotus 1-2-3   & everything                 & topological, found on the fly \\
\typ{make}    & targets older than their inputs & topological \\
Excel         & dirty cells and their dependents & self-repairing chain (``restarting'') \\
\bottomrule
\end{tabularx}
\end{center}

The same idea, recomputing only what depends on what changed, runs through
\emph{incremental computation} research (self-adjusting computation,
Adapton), reactive notebooks such as Observable, and the \emph{signals} of
today's web frameworks: Solid, Preact, Angular, and a 2024 proposal to add
signals to JavaScript itself. A signal is a cell, a computed signal is a
formula, and the framework is a recalculation engine.

So, for now, the history ends with the spreadsheet's recalculation engine as
a general model for how interactive software keeps itself up to date.

% ======================================================================
\section{Kalkylark's engines}

Kalkylark implements the first two stages, and the code is meant to be
read.

\subsection{The sweep}

The default, \typ{/GOC} (or \typ{/GOR}), is VisiCalc's single pass.
Formulas are stored as text and evaluated by a recursive-descent parser
that evaluates as it parses, strictly left to right, reading each
referenced cell's last computed value.

\subsection{Natural order}

\typ{/GON} switches to natural order. It takes three steps:

\begin{enumerate}[label=\arabic*.]
\item \textbf{Find the edges.} \typ{formula\_refs()} runs the same parser
  in a mode where every cell a formula reads is reported. A range reports
  all of its cells. \typ{@LOOKUP} also reports the column (or row) its
  result is taken from. Which cell that is depends on the data, so this is
  a conservative stand-in for a dependency that is only known at run time,
  the same problem Excel's \typ{INDIRECT} has.
\item \textbf{Order them.} Rather than 1-2-3's recursive walk, Kalkylark
  uses \emph{Kahn's algorithm}. Every formula that reads no other formula
  goes in a queue, and each dependent is released when its last precedent
  has been computed. There is no recursion, so even a chain through all
  16\,002 cells cannot overflow the stack.
\item \textbf{Handle cycles.} Formulas that are never released lie on a
  cycle, or downstream of one. They get one VisiCalc-style pass and raise
  the \textsf{CIRC} flag, so a deliberately circular model converges with
  repeated \typ{!}.
\end{enumerate}

Minimal recalculation, the Excel stage, would add a list of dependents per
cell and dirty marking. On a 63~by~254 sheet with modern hardware it is
rarely needed, but it would be an instructive next step. The edge-finding
code it needs already exists.

% ======================================================================
\section{Rights: VisiCalc today}

A reconstruction raises the question of who owns the original, and what
that means for a new implementation. This section summarises what we found,
as of November 2026. It is a summary, not legal advice.

\subsection{Who owns what}

\begin{description}[style=nextline, leftmargin=1.5em]
\item[The code.] VisiCalc was written by Software Arts (Bricklin and
  Frankston) and published by Personal Software, later VisiCorp. After a
  dispute between the two companies, settled in 1984, Lotus Development
  Corporation acquired Software Arts and the VisiCalc rights in 1985 and
  stopped selling the program. A later court ruling confirmed that Lotus
  had acquired all rights. IBM acquired Lotus in 1995. A US copyright in a
  work published in 1979 by a company runs for 95 years, so the program's
  code remains protected for decades yet.
\item[The name.] Contemporary reports of the 1984 settlement say Software
  Arts kept the \emph{VisiCalc} name and VisiCorp the ``Visi'' prefix. We
  found no current registration either way. Kalkylark has its own name, so
  the question does not arise.
\item[The idea.] VisiCalc was never patented. Bricklin has said he was
  advised that a patent was unlikely to be granted for software at the time.
  The idea of the electronic spreadsheet was free from the start.
\end{description}

\subsection{The licence of the original program}

Dan Bricklin's web site offers the original IBM PC version, \texttt{VC.COM},
under a licence from Lotus, which is named on the page as the owner. We
read the page on 19 November 2026. In short:

\begin{itemize}
\item use is permitted for \emph{personal, non-commercial} purposes, with
  copies only to support that use and with the copyright notice kept;
\item copying, modifying or distributing the program beyond that is not
  permitted, nor is it permitted to ``reverse assemble, reverse compile, or
  otherwise translate'' it, except where the law gives that right
  regardless of contract;
\item the original reference card is covered by the same licence. It may
  be printed for personal use but not distributed.
\end{itemize}

The licence covers that executable and that card. It says nothing about
clones, file formats or the spreadsheet concept.

\subsection{Why a reconstruction is on safe ground}

Copyright protects a program's \emph{expression}, its code, and not its
ideas or the way it works. Courts on both sides of the Atlantic have drawn
the line clearly:

\begin{itemize}
\item \textbf{European Union.} The Software Directive (2009/24/EC) states
  that ideas and principles underlying a program are not protected. In
  \emph{SAS Institute v World Programming} (case C-406/10, 2012), the Court
  of Justice held that a program's functionality, its programming language
  and the format of its data files are not protected by copyright. Copying
  the behaviour of a program you have studied, without copying its code, is
  lawful.
\item \textbf{United States.} In \emph{Lotus v.\ Borland} (1st Cir.\ 1995,
  affirmed by an evenly divided Supreme Court in 1996), Lotus 1-2-3's menu
  command hierarchy, the very kind of \typ{/} command tree that VisiCalc
  invented, was held to be an uncopyrightable \emph{method of operation}.
  An earlier district court decision, \emph{Lotus v.\ Paperback Software}
  (1990), had gone against a 1-2-3 clone on ``look and feel''. Borland
  largely displaced it, but it shows where the risk once was.
\end{itemize}

Kalkylark therefore follows a simple rule: write everything from scratch,
recreate the behaviour, and use none of the original's code, artwork, name
or documentation. The original program and its reference card may be used,
under Lotus's licence, to check behaviour, but they are never disassembled,
and they are not distributed with Kalkylark.

% ======================================================================
\section*{Sources}
\addcontentsline{toc}{section}{Sources}

\begin{small}
\begin{itemize}[leftmargin=1.2em, itemsep=2pt]
\item Microsoft Learn, \emph{Excel Recalculation}; \emph{Multithreaded
  Recalculation in Excel}; \emph{Improving calculation performance}.
\item A.~Mokhov, N.~Mitchell, S.~Peyton Jones, \emph{Build systems à la
  carte}, Proc.\ ACM Program.\ Lang.\ 2 (ICFP), 2018.
\item D.~Bricklin, \emph{VisiCalc License} and history pages,
  \url{http://www.bricklin.com/history/} (read 19 November 2026).
\item Computer History Museum, oral history of Dan Bricklin and Peter
  Fylstra (catalogue 102701923).
\item ``Software Arts Inc., VisiCorp settle suits over VisiCalc program'',
  1984 news report.
\item US Patent 4,398,249 (Pardo and Landau), granted 1983.
\item Court of Justice of the EU, case C-406/10, \emph{SAS Institute Inc.\
  v World Programming Ltd}, 2 May 2012. Directive 2009/24/EC, Articles 1,
  5 and 6.
\item \emph{Lotus Development Corp.\ v.\ Borland International, Inc.}, 49
  F.3d 807 (1st Cir.\ 1995), aff'd by an equally divided Court, 516 U.S.
  233 (1996). \emph{Lotus Development Corp.\ v.\ Paperback Software
  International}, 740 F.\ Supp.\ 37 (D.\ Mass.\ 1990).
\end{itemize}
\end{small}

\vfill
{\small This document is licensed CC BY 4.0. VisiCalc is a trademark of
its respective owner. Kalkylark is not affiliated with Lotus, IBM,
Software Arts or VisiCorp.}

\end{document}
